\chapter{可视化界面设计与实现}

\section{设计目标}

除模型训练与离线测试外,本文还实现了一个面向演示与验证的可视化界面,用于展示细粒度分类系统在实际交互场景中的运行方式。增加可视化界面的目的主要有三点。第一,将命令行环境中的模型推理流程转化为可直接操作的页面,降低实验展示门槛。第二,把图像选择、模型运行和结果输出组织为连续步骤,便于观察系统从输入到预测的完整链路。第三,通过页面结果对比,更直观地体现本文方法与原始 CLIP 在细粒度分类任务上的差异。

本界面采用 Streamlit 实现。选择这一方案的原因在于,它能够较快完成原型搭建,并且适合将 Python 推理脚本直接封装为网页交互逻辑。对于毕业设计场景而言,这种实现方式既能满足功能展示需求,又不需要额外构建复杂的前后端分离系统。

\section{界面整体布局}

界面整体采用单页式布局,将功能区域按照用户使用顺序自上而下组织。页面顶部用于说明系统名称、任务目标和使用方式,中部为图像输入与参数选择区域,下部用于展示模型预测结果和对比信息。这样的布局能够减少操作跳转,使用户在一次页面浏览中完成上传、推理和查看结果的完整流程。

图~\ref{fig:streamlit_layout} 展示了可视化界面的整体布局。从页面结构上看,界面重点突出了输入区与结果区的对应关系,便于在演示时快速说明系统包含哪些核心功能模块。

\begin{figure}[H]
\centering

\includegraphics[width=0.9\textwidth]{layout_of_the_streamlit_demo.png}
\caption{Streamlit 演示界面的整体布局}
\label{fig:streamlit_layout}
\end{figure}

\section{图像选择与推理流程}

在具体交互流程上,用户首先需要选择待识别图像,然后触发模型推理,最后查看分类输出。为了适配细粒度分类任务,界面中保留了图片预览区域,使用户在运行前能够先确认输入内容是否正确。这一点对于宠物品种识别尤其重要,因为不同类别之间差异较小,错误选择图片会直接影响结果判断。

图~\ref{fig:streamlit_select} 给出了界面中的图像选择与运行区域。可以看到,系统将待测图片选择、运行按钮和结果触发逻辑集中在同一位置,这样能够减少操作分散带来的干扰。对于展示型系统而言,这种设计比完全依赖脚本参数输入更直观,也更符合非开发者的使用习惯。

\begin{figure}[H]
\centering
\includegraphics[width=0.9\textwidth]{select_picture_to_run_model.png}
\caption{待测图片选择与模型运行区域}
\label{fig:streamlit_select}
\end{figure}

在实现上,界面前端负责接收图片输入并调用后端推理逻辑,后端则加载已经训练完成的 DynamicPromptTrainer 模型,输出类别预测结果及其置信信息。整个过程不重新训练模型,而是直接复用已有 checkpoint,因此页面响应的重点在于推理效率与展示清晰度,而不是训练过程管理。

\section{结果展示与方法对比}

可视化界面不仅用于给出单个预测结果,也用于展示本文方法与基线方法之间的差异。相比只输出一个类别标签,页面进一步给出不同方法的分类结果对比,可以更直观地说明动态提示优化在细粒度识别中的实际效果。对于毕业设计答辩或项目展示场景,这种结果呈现方式比单纯引用表格更容易让观察者理解模型改进带来的收益。

图~\ref{fig:web_result_compare} 展示了网页端的结果对比示例。其中,本文方法与原始 CLIP 在同一输入上的输出被放在统一页面中进行展示,有助于直接观察两者在预测类别和表现稳定性上的区别。这样的设计使可视化界面不仅是一个“能运行模型”的工具,也成为实验结果解释与方法展示的一部分。

\begin{figure}[H]
\centering
\includegraphics[width=0.9\textwidth]{web_result_ours_vs_clip.png}
\caption{网页端本文方法与原始 CLIP 的结果对比}
\label{fig:web_result_compare}
\end{figure}

\section{本章小结}

本章介绍了本文实现的可视化界面工作。基于 Streamlit 的单页式界面将图像选择、模型推理和结果展示整合到统一页面中,使细粒度分类系统具备了更直观的演示与交互能力。通过加入本文方法与原始 CLIP 的结果对比,该界面不仅提升了系统展示效果,也为实验结论提供了更易理解的可视化支撑。
